\documentclass[12pt]{article}
\usepackage[utf8]{inputenc}
\usepackage{amsmath}
\usepackage{amssymb}
\usepackage{geometry}
\usepackage{hyperref}

\geometry{a4paper, margin=1in}

\title{\texttt{Orbit} \\ \large Continuous Navigation via Actor-Critic \\ Methods using Neural Function Approximation}
\author{Stochastic Batman}
\date{March 16-17, 2026}

\begin{document}
	
	\maketitle
	
	\section*{Introduction}
	This document outlines the thinking process and systematic design for \href{https://github.com/Stochastic-Batman/Orbit}{\texttt{Orbit}}. It serves as a bridge between the theoretical foundations established in \textit{Mathematical Foundations of Reinforcement Learning} by Shiyu Zhao and the practical implementation of neural function approximation in a continuous coordinate space.
	
	\section{Environment Specification (\texttt{env.rs})}
	
	The environment is modeled as a Markov Decision Process (MDP) where the agent must learn to navigate toward a target in a bounded 2D plane.
	
	\subsection{State Space $\mathcal{S}$}
	
	The state space is defined over a four-dimensional hypercube. To facilitate stable training in neural networks, all coordinates are normalized to the range $[-1, 1]$. I define the state vector $s \in \mathcal{S}$ as:
	$$ s \doteq [x, y, x_{target}, y_{target}]^\top \in [-1, 1]^4 $$ 
	To maintain consistency in subsequent notation, I denote the components of any state $s$ as:
	\begin{itemize}
		\item $s[0] = x$ (Agent's horizontal position) and $s[1] = y$ (Agent's vertical position)
		\item $s[2] = x_{target}$ (Target's horizontal position) and $s[3] = y_{target}$ (Target's vertical position)
	\end{itemize}
	\textbf{Notation Note:} I use the shorthand $s[i, \dots, k]$ to represent a sub-vector consisting of elements from index $i$ to $k$ inclusive. For example, $s[0, 1]$ refers to the agent's coordinate pair $[s[0], s[1]]^\top$.
	
	\subsection{Action Space $\mathcal{A}$ and Transition Dynamics}
	
	I employ a discrete action space $\mathcal{A} = \{a_1, a_2, a_3, a_4, a_5\}$. Each action corresponds to a displacement vector $\Delta a \in \mathbb{R}^2$ with a step size magnitude $\Delta = 0.05$. The physical directions follow the clockwise convention:
	
	$$ \mathcal{A} = \{ 
	\overset{\text{North}}{a_1}, 
	\overset{\text{East}}{a_2}, 
	\overset{\text{South}}{a_3}, 
	\overset{\text{West}}{a_4}, 
	\overset{\text{Stay}}{a_5} 
	\} $$
	
	The state transition probability $p(s' | s, a)$ is deterministic. The next state $s'$ is calculated by applying the displacement to the agent's coordinates $s[0, 1]$ and clipping the result to the boundary: 
	$$ s'[i] = \text{clip}(s[i] + \Delta a_i, -1, 1), \, i \in \{0, 1\} $$ 
	The target coordinates remain stationary: $s'[2, 3] = s[2, 3]$.
	
	\subsection{Reward Function $r(s, a)$}
	The reward signal provides dense feedback based on the Euclidean distance between the agent and the target. I define the distance function between any two points $p_1, p_2$ as:
	$$ \text{dist}(p_1, p_2) \doteq \sqrt{ (p_1[0] - p_2[0])^2 + (p_1[1] - p_2[1])^2 } $$
	
	The reward function is formulated to ensure the boundary penalty is strictly lower than any possible reward achievable within the valid state space. In the square $[-1, 1]^2$, the maximum possible distance (the diagonal) is $\sqrt{2^2 + 2^2} = \sqrt{8} \approx 2.8284$. 
	
	I define the penalty as $-2.85$. This extra $\approx -0.02$ is to ensure that a boundary collision is always perceived by the Critic as a "worse" state than being at the furthest possible valid coordinate from the target.
	$$ 
	r(s, a) = 
	\begin{cases} 
		-2.85 & \text{if } s[0,1] + \Delta a \notin [-1, 1]^2 \text{ (Boundary Crossing)} \\
		-\text{dist}(s'[0,1], s'[2,3]) & \text{otherwise}
	\end{cases}
	$$
	
	\subsection{Termination and Episodes}
	The task is episodic with the following conditions:
	\begin{itemize}
		\item \textbf{Success Condition:} The episode terminates when $\text{dist}(s[0,1], s[2,3]) < 0.05$.
		\item \textbf{Truncation Condition:} An episode is capped at $T_{max} = 250$ steps.
		\item \textbf{Terminal Value:} Following the book's convention, the value of the terminal state $v(s_{terminal})$ is defined as $0$.
	\end{itemize}
	
	\section{Reinforcement Learning Framework}
	
	Following the Actor-Critic framework (Zhao, Chapter 10), I decouple the neural architectures from the training logic to ensure modularity.
	
	\subsection{Neural Architectures (\texttt{model.rs})}
	I define the "blueprints" for the function approximators using the \texttt{dfdx} type-system. Both networks receive the state $s \in [-1, 1]^4$.
	\begin{itemize}
		\item \textbf{Actor $\pi_\theta$}: A 3-layer MLP mapping states to 5 action logits.
		\item \textbf{Critic $v_w$}: A 3-layer MLP mapping states to a scalar value estimate.
	\end{itemize}
	
	\subsection{The Actor-Critic Agent (\texttt{agent.rs})}
	The \texttt{OrbitAgent} manages the interaction between the networks and the environment's feedback loop.
	
	\subsubsection{Temporal Difference and Advantage}
	We utilize the TD error $\delta_t$ as an estimate of the advantage $A(s,a)$:
	$$ \delta_t = R_{t+1} + \gamma \hat{v}(S_{t+1}, w) - \hat{v}(S_t, w) $$
	The agent handles termination by setting $\hat{v}(S_{t+1}, w) = 0$ if \texttt{is\_done} is true. For \texttt{is\_truncated}, we bootstrap to preserve the infinite-horizon MDP assumption.
	
	\subsubsection{Objective Functions}
	The agent performs two distinct gradient updates per step:
	\begin{enumerate}
		\item \textbf{Value Update}: Minimizes $L(w) = \frac{1}{2} \delta_t^2$ via backpropagation through the Critic.
		\item \textbf{Policy Update}: Minimizes the weighted negative log-likelihood with entropy regularization:
		$$ L(\theta) = - \ln \pi(A_t | S_t, \theta) \delta_{t, \text{stop\_grad}} - \beta H(\pi) $$
		where $\delta_{t, \text{stop\_grad}}$ indicates that the TD error is treated as a constant scalar during the actor's backpropagation. "$-$" before $\ln$ is because \texttt{dfdx} minimizes instead of maximizing.
	\end{enumerate}
	
	\subsection{Optimization and Hyperparameters}
	The parameters $\theta$ and $w$ are updated using the Adam optimizer. To ensure stable convergence, I employ a two-time-scale update rule where the Critic typically learns faster than the Actor to provide a stable baseline for the advantage signal:
	\begin{itemize}
		\item \textbf{Discount Factor ($\gamma$)}: $0.99$ (Prioritizing long-term navigation).
		\item \textbf{Learning Rates}: $\alpha_\theta = 1 \times 10^{-4}$ (Actor), $\alpha_w = 1 \times 10^{-3}$ (Critic).
		\item \textbf{Entropy Regularization}: To prevent premature convergence to a sub-optimal deterministic policy, I may add an entropy term $H(\pi)$ to the Actor loss: $L_{total}(\theta) = L(\theta) - \beta H(\pi(s, \theta))$ where $\beta = 0.01$ .
	\end{itemize}
	
	\subsection{Softmax Policy Derivation}
	For the discrete action set $\{a_1, \dots, a_5\}$, the probability of selecting action $a_i$ given logits $z$ from the Actor is:
	$$ \pi(a_i | s, \theta) = \frac{e^{z_i(s, \theta)}}{\sum_{j=1}^5 e^{z_j(s, \theta)}} $$
	
\end{document}